\chapter{Normed Vector Spaces}\label{ch:b2:nvs}

When the metric space is a vector space and the distance comes from
a norm, topology and linear algebra begin to interact: linear maps
are continuous exactly when they are bounded on the unit ball,
finite dimension forces all norms to agree, and completeness turns
absolutely convergent series into convergent ones. The
finite/infinite dimension divide --- crystallized in Riesz's theorem
--- is the deepest lesson of the chapter.

Throughout, $E, F$ are vector spaces over $K = \R$ or $\C$.

\section{Norms}

\begin{definition}\label{def:b2:nvs:norm}
A \emph{norm}\index{norm} on $E$ is a map $\norm{\,\cdot\,} \colon E
\to \R_+$ with, for all $x, y \in E$, $\lambda \in K$:
\[
\norm x = 0 \iff x = 0,
\qquad
\norm{\lambda x} = \abs\lambda\,\norm x,
\qquad
\norm{x + y} \leq \norm x + \norm y .
\]
Then $d(x, y) = \norm{x - y}$ is a distance, and all of
\cref{ch:b2:metric} applies. The reverse triangle inequality
$\bigl|\norm x - \norm y\bigr| \leq \norm{x - y}$ makes the norm
itself $1$-Lipschitz; addition and scalar multiplication are
continuous (estimates $\norm{(x + y) - (x' + y')} \leq \norm{x -
x'} + \norm{y - y'}$, etc.).
\end{definition}

\begin{example}\label{ex:b2:nvs:examples}
On $K^n$:
\[
\norm{x}_1 = \sum_i \abs{x_i},
\qquad
\norm{x}_2 = \Bigl(\sum_i \abs{x_i}^2\Bigr)^{1/2},
\qquad
\norm{x}_\infty = \max_i \abs{x_i}
\]
($\norm\cdot_2$ is a norm by Cauchy--Schwarz, Year 1 volume). On
$C(\intcc{a}{b})$:
\[
\norm f_\infty = \sup \abs f,
\qquad
\norm f_1 = \int_a^b \abs f,
\qquad
\norm f_2 = \Bigl(\int_a^b \abs f^2\Bigr)^{1/2},
\]
the last two being norms thanks to strict positivity of the
integral and integral Cauchy--Schwarz (Year 1 volume). On matrices:
any norm on $\mathcal M_n(K) \simeq K^{n^2}$; the operator norms
below are the structurally important ones.
\end{example}

\begin{definition}[Equivalent norms]\label{def:b2:nvs:equivalent}
Two norms $N_1, N_2$ on $E$ are \emph{equivalent}\index{equivalent
norms} when there are constants $c, C > 0$ with
\[
c\,N_1 \leq N_2 \leq C\, N_1 .
\]
Equivalent norms have the same open sets, the same convergent and
Cauchy sequences, the same compact and complete subsets: the same
analysis.
\end{definition}

\begin{example}[Non-equivalence in infinite dimension]\label{ex:b2:nvs:nonequivalent}
On $C(\intcc{0}{1})$: $\norm f_1 \leq \norm f_\infty$ always, but
no reverse bound holds: $f_n(x) = x^n$ has $\norm{f_n}_\infty = 1$
and $\norm{f_n}_1 = \frac{1}{n+1} \to 0$. So $f_n \to 0$ for
$\norm\cdot_1$ but not for $\norm\cdot_\infty$: the two norms
disagree about convergence itself.
\end{example}

\begin{example}[Explicit constants in dimension $n$]\label{ex:b2:nvs:constants}
On $K^n$ the three classical norms are equivalent with sharp
constants:
\[
\norm x_\infty \leq \norm x_2 \leq \norm x_1
\leq \sqrt n\,\norm x_2 \leq n\,\norm x_\infty ,
\]
the middle bound $\norm x_1 \leq \sqrt n\norm x_2$ coming from
Cauchy--Schwarz against the all-ones vector. Extremal vectors:
$e_1$ makes the first two inequalities equalities, $(1, 1,
\dots, 1)$ the last two. The dimension $n$ sits visibly in the
constants --- the quantitative seed of the failure in infinite
dimension: as $n \to \infty$ no uniform constant survives,
which is exactly what \cref{ex:b2:nvs:nonequivalent} exhibits on
function spaces.
\end{example}

\begin{omfigure}
\begin{tikzpicture}[scale=1.5]
  \draw[->] (-1.45,0) -- (1.5,0) node[below, font=\small] {$x$};
  \draw[->] (0,-1.45) -- (0,1.5) node[left, font=\small] {$y$};
  \draw[omProp, very thick]
    (1,1) -- (-1,1) -- (-1,-1) -- (1,-1) -- cycle;
  \draw[omThm, very thick] (0,0) circle (1);
  \draw[omDef, very thick]
    (1,0) -- (0,1) -- (-1,0) -- (0,-1) -- cycle;
  \node[omDef, font=\small] at (0.36,0.28)
    {$\norm\cdot_1$};
  \node[omThm, font=\small] at (0.55,0.92)
    {$\norm\cdot_2$};
  \node[omProp, font=\small] at (1.24,1.13)
    {$\norm\cdot_\infty$};
\end{tikzpicture}

{\small The unit balls of the three classical norms of $\R^2$,
nested as the inequalities of \cref{ex:b2:nvs:constants}
dictate: smaller ball, larger norm. Roundness matters: the flat
sides of the diamond and the square are exactly the failures of
strict convexity exploited in \cref{ch:b2:realfun}'s weekend
problem and in this chapter's (question 4).}
\end{omfigure}

\section{Continuous linear maps}

\begin{theorem}[Characterization]\label{thm:b2:nvs:continuouslinear}
For a linear map $u \colon E \to F$ between normed spaces, the
following are equivalent:
\begin{enumerate}
  \item $u$ is continuous;
  \item $u$ is continuous at $0$;
  \item $u$ is bounded on the closed unit ball: $\sup_{\norm x \leq
        1} \norm{u(x)} < \infty$;
  \item there is $C \geq 0$ with $\norm{u(x)} \leq C \norm x$ for
        all $x$;
  \item $u$ is Lipschitz.
\end{enumerate}
The least such $C$ is the \emph{operator norm}\index{operator norm}
$\vertiii{u} = \sup_{\norm x \leq 1}\norm{u(x)} = \sup_{x \neq 0}
\frac{\norm{u(x)}}{\norm x}$; it makes the space $\mathcal{L}_c(E,
F)$ of continuous linear maps a normed space, with
\[
\vertiii{v \circ u} \leq \vertiii v\, \vertiii u .
\]
\end{theorem}

\begin{proof}
(1 $\Rightarrow$ 2) trivial. (2 $\Rightarrow$ 3): continuity at $0$
with $\varepsilon = 1$ gives $\delta$ with $\norm x \leq \delta
\Rightarrow \norm{u(x)} \leq 1$; homogeneity then scales any $x$
with $\norm x \leq 1$ down into that ball and back:
\[
\norm{u(x)} = \frac1\delta\,\norm{u(\delta x)} \leq
\frac1\delta ,
\]
since $\norm{\delta x} \leq \delta$. (3 $\Rightarrow$ 4):
for $x \neq 0$, apply the bound to $\frac{x}{\norm x}$. (4
$\Rightarrow$ 5): $\norm{u(x) - u(y)} = \norm{u(x - y)} \leq
C\norm{x - y}$. (5 $\Rightarrow$ 1) known.

Norm axioms for $\vertiii\cdot$: homogeneity and separation are
clear ($\vertiii u = 0$ forces $u = 0$ on the ball, hence
everywhere); triangle inequality from $\norm{(u + v)(x)} \leq
\norm{u(x)} + \norm{v(x)}$. Submultiplicativity: $\norm{v(u(x))}
\leq \vertiii v\,\norm{u(x)} \leq \vertiii v \vertiii u \norm x$.
\end{proof}

\begin{example}\label{ex:b2:nvs:operatornorms}
On $\bigl(C(\intcc{0}{1}), \norm\cdot_\infty\bigr)$: evaluation
$f \mapsto f(0)$ has operator norm $1$; integration $f \mapsto
\int_0^1 f$ has norm $1$; the map $f \mapsto \int_0^1 t f(t)\dd t$
has norm $\int_0^1 t\,\dd t = \frac12$ (upper bound by the triangle
inequality for integrals; attained at $f \equiv 1$). But
differentiation, from $(C^1, \norm\cdot_\infty)$ to $(C^0,
\norm\cdot_\infty)$, is \emph{not} continuous: $\norm{\sin(nx)}_\infty
= 1$ while the derivative has sup norm $n$. Linear does not imply
continuous in infinite dimension.
\end{example}

\begin{example}[Two norms, two verdicts on one sequence]\label{ex:b2:nvs:sqrtnxn}
On $C(\intcc01)$, let $g_n(x) = \sqrt{n}\,x^n$. Then
\[
\norm{g_n}_1 = \frac{\sqrt n}{n + 1} \longrightarrow 0,
\qquad
\norm{g_n}_2^2 = \frac{n}{2n + 1} \longrightarrow \frac12,
\qquad
\norm{g_n}_\infty = \sqrt n \longrightarrow \infty :
\]
one sequence, three norms, three behaviours --- convergence to
zero, no convergence (norms stabilize at $\frac1{\sqrt2}$ but
the pointwise limit is $0$), and explosion. Mass concentrating
near $x = 1$ is invisible to $\norm\cdot_1$, half-visible to
$\norm\cdot_2$, dominant for $\norm\cdot_\infty$. In infinite
dimension, ``does it converge?'' is not a question about a
sequence: it is a question about a sequence \emph{and} a norm.
\end{example}

\begin{method}[Computing an operator norm]\label{met:b2:nvs:opnorm}
Always in two moves. \emph{Upper bound:} estimate
$\norm{u(x)}$ by $C\norm x$ using triangle inequalities,
Cauchy--Schwarz, or integral bounds --- this proves $\vertiii u
\leq C$. \emph{Witness:} exhibit either a specific $x_0 \neq 0$
with $\norm{u(x_0)} = C\norm{x_0}$ (the bound is attained), or a
sequence of unit vectors $x_n$ with $\norm{u(x_n)} \to C$ (the
bound is approached). Both moves are mandatory: an upper bound
alone gives only $\vertiii u \leq C$, a witness alone only
$\vertiii u \geq C$. In infinite dimension the witness may have
to be a sequence --- the supremum need not be attained
(\cref{exo:b2:nvs:8}).
\end{method}

\begin{example}[Diagonal operators see every norm alike]\label{ex:b2:nvs:diagonalnorm}
For $D = \operatorname{diag}(d_1, \dots, d_n)$ on $K^n$ with
\emph{any} of the norms $\norm\cdot_1, \norm\cdot_2,
\norm\cdot_\infty$: from $\abs{d_ix_i} \leq
\bigl(\max_j\abs{d_j}\bigr)\abs{x_i}$ coordinatewise,
$\norm{Dx} \leq \max_j\abs{d_j}\,\norm x$; and $x = e_{j_0}$
(a maximizing index) attains it. So $\vertiii D =
\max_j\abs{d_j}$ in all three cases: for diagonal maps, all
reasonable norms tell the same story, the largest stretch
factor. Everything difficult about operator norms is about
\emph{non}-diagonal behaviour --- which is why the adapted
norms of this chapter's weekend problem (question 22) work by
forcing a matrix to become diagonal first.
\end{example}

\begin{example}[Column sums: the $1$-norm twin of \cref{exo:b2:nvs:4}]\label{ex:b2:nvs:columnsums}
On $(\R^n, \norm\cdot_1)$, the operator norm of a matrix $A$ is
the largest absolute \emph{column} sum. Run the method: for
$\norm x_1 \leq 1$,
\[
\norm{Ax}_1 = \sum_i\Bigl|\sum_j a_{ij}x_j\Bigr|
\leq \sum_j \abs{x_j}\sum_i\abs{a_{ij}}
\leq \Bigl(\max_j\sum_i\abs{a_{ij}}\Bigr)\norm x_1 ,
\]
and the bound is attained at $x = e_{j_0}$ for a maximizing
column $j_0$ --- the tidiest witness imaginable. So for $A =
\begin{psmallmatrix}1 & -2\\ 3 & 1\end{psmallmatrix}$:
$\vertiii A_1 = \max(1 + 3,\ 2 + 1) = 4$, while
$\vertiii A_\infty = 4$ too (rows) --- a coincidence here, not
a law: transpose the matrix entries asymmetrically and the two
norms part company. Rows for $\norm\cdot_\infty$, columns for
$\norm\cdot_1$: the mnemonic is that each norm's unit vectors
(sign patterns, resp.\ basis vectors) pick out the matching
sums.
\end{example}

\begin{proposition}[Bilinear maps]\label{prop:b2:nvs:bilinear}
A bilinear map $b \colon E \times F \to G$ is continuous iff
$\norm{b(x,y)} \leq C\norm x\,\norm y$ for some $C$; then it is
Lipschitz on bounded sets. (Same proof pattern; the product
$(u, v) \mapsto v \circ u$ and matrix multiplication are the key
examples.)
\end{proposition}

\begin{proof}
If the bound holds:
\[
b(x,y) - b(x_0,y_0) = b(x - x_0,\, y) + b(x_0,\, y - y_0),
\]
so $\norm{b(x,y) - b(x_0,y_0)} \leq C\norm{x - x_0}\norm y +
C\norm{x_0}\norm{y - y_0}$: continuity at $(x_0, y_0)$, and a
Lipschitz bound where $\norm x, \norm y \leq R$. Conversely,
continuity at $(0,0)$ gives $\delta$ with $\norm{b(x,y)} \leq 1$ on
$\norm x, \norm y \leq \delta$; scale both variables.
\end{proof}

\section{Finite dimension}

\begin{theorem}[Equivalence of norms in finite dimension]\label{thm:b2:nvs:finitedim}
On a finite-dimensional space, \emph{all norms are
equivalent}.\index{equivalence of norms} Consequently, in finite
dimension: convergence, openness, compactness, completeness are
norm-independent notions; compact $=$ closed and bounded;
the space is complete; and every linear (or multilinear) map
\emph{from} a finite-dimensional space is continuous.
\end{theorem}

\begin{proof}
Fix a basis and identify $E \simeq K^n$; it suffices to compare any
norm $N$ with $\norm\cdot_\infty$.

\emph{One direction is algebra:} $N(x) = N(\sum x_i e_i) \leq
\sum \abs{x_i} N(e_i) \leq C \norm x_\infty$ with $C = \sum N(e_i)$.
This also shows $N$ is continuous on $(K^n, \norm\cdot_\infty)$
(it is $C$-Lipschitz: $\abs{N(x) - N(y)} \leq N(x - y)$).

\emph{The other is topology:} the unit sphere $S = \{x :
\norm{x}_\infty = 1\}$ is closed and bounded in $(K^n,
\norm\cdot_\infty)$, hence compact
(\cref{thm:b2:metric:compactprops}~(2), valid for $\C^n \simeq
\R^{2n}$). The continuous function $N$ attains its minimum $c$ on
$S$; $c > 0$ since $N$ vanishes only at $0 \notin S$. Homogeneity
spreads the bound: $N(x) \geq c \norm{x}_\infty$ for all $x$.

Consequences: all statements reduce to $(K^n,
\norm\cdot_\infty)$, where they are known
(\cref{thm:b2:metric:rncomplete},
\cref{thm:b2:metric:compactprops}); a linear $u$ from
finite-dimensional $E$ satisfies $\norm{u(x)} \leq \sum\abs{x_i}
\norm{u(e_i)} \leq C'\norm{x}_\infty$: bound (4) of
\cref{thm:b2:nvs:continuouslinear}.
\end{proof}

\begin{corollary}\label{cor:b2:nvs:closedsubspace}
A finite-dimensional subspace of any normed space is closed.
\end{corollary}

\begin{proof}
It is complete for the induced norm
(\cref{thm:b2:nvs:finitedim}), and complete subsets are closed
(\cref{def:b2:metric:complete}).
\end{proof}

\begin{example}[A best approximation computed by symmetry]\label{ex:b2:nvs:bestaffine}
In $\bigl(C(\intcc{-1}{1}), \norm\cdot_\infty\bigr)$, how far is
$f(x) = \abs x$ from the (closed, two-dimensional) subspace of
affine functions $a + bx$? By symmetry, replacing $a + bx$ by
$a - bx$ leaves $\norm{f - (a \pm bx)}_\infty$ unchanged, and
the midpoint $a$ does at least as well (triangle inequality on
the average): it suffices to consider constants. For a constant
$a$:
\[
\norm{\abs x - a}_\infty = \max\,(1 - a,\ a)
\geq \frac12 ,
\]
minimized at $a = \frac12$: the distance is $\frac12$, attained
by the constant $\frac12$. Note the error curve $\abs x -
\frac12$: it reaches $\pm\frac12$ alternately at $x = -1, 0,
1$ --- three extrema of alternating sign for a best
approximation from a two-parameter family. That
\emph{equioscillation} pattern is no accident; it is the
signature of optimality that this chapter's weekend problem
turns into Chebyshev's theorem.
\end{example}

\begin{example}[Closed versus dense subspaces]\label{ex:b2:nvs:closeddense}
In $E = \bigl(C(\intcc{0}{1}), \norm\cdot_\infty\bigr)$: each
$\R_n[X]$ (polynomials of degree $\leq n$, restricted to
$\intcc01$) is a finite-dimensional, hence \emph{closed}, subspace
--- a uniform limit of polynomials of degree $\leq n$ is one. But
the union $\R[X]$ of all of them is \emph{dense} in $E$ (the
Weierstrass approximation theorem, proved in
\cref{ch:b2:funcseq}), and dense proper subspaces are as
non-closed as can be. The moral: closedness of subspaces is a
finite-dimensional privilege; stacking closed floors can build a
dense skyscraper.
\end{example}

\begin{theorem}[Riesz]\label{thm:b2:nvs:riesz}
The closed unit ball of a normed space $E$ is compact \emph{if and
only if} $\dim E < \infty$.\index{Riesz's theorem}
\end{theorem}

\begin{proof}
Finite dimension: closed and bounded suffices
(\cref{thm:b2:nvs:finitedim}).

Conversely, suppose $\dim E = \infty$. \emph{Riesz's lemma:} for
every proper closed subspace $F \subsetneq E$ and $\varepsilon \in
\intoo{0}{1}$, there is a unit vector $x$ with $d(x, F) \geq 1 -
\varepsilon$. Proof: pick $y \notin F$, let $\delta = d(y, F) > 0$
($F$ closed), choose $f \in F$ with $\norm{y - f} \leq
\frac{\delta}{1 - \varepsilon}$, and set $x = \frac{y - f}{\norm{y -
f}}$: for any $g \in F$,
\[
\norm{x - g} = \frac{\norm{y - (f + \norm{y-f}\,g)}}{\norm{y - f}}
\geq \frac{\delta}{\norm{y-f}} \geq 1 - \varepsilon ,
\]
the numerator being a distance from $y$ to a point of $F$.

Now build unit vectors $x_1, x_2, \dots$ inductively: $F_k =
\operatorname{Vect}(x_1, \dots, x_k)$ is finite-dimensional, hence
closed (\cref{cor:b2:nvs:closedsubspace}) and proper; Riesz's lemma
with $\varepsilon = \frac12$ provides a unit $x_{k+1}$ with $d(x_{k+1},
F_k) \geq \frac12$. The sequence satisfies $\norm{x_p - x_q} \geq
\frac12$ for $p \neq q$: no convergent subsequence --- the unit ball
is not compact.
\end{proof}

\begin{example}[Riesz as a dimension detector]\label{ex:b2:nvs:rieszdetector}
Is $C(\intcc01)$ finite-dimensional? Riesz answers without
exhibiting any explicit infinite free family: the sequence
$f_n(x) = x^n$ lies in the closed unit ball and satisfies, for
$m > n$, $\norm{f_n - f_m}_\infty \geq f_n(x_0) - f_m(x_0) > 0$
at suitable points --- quantified cleanly in this chapter's
weekend problem (question 16), where a subsequence stays at
mutual distance $\geq \frac14$. No convergent subsequence, so
the ball is not compact, so $\dim C(\intcc01) = \infty$ by
\cref{thm:b2:nvs:riesz}. Compactness of the unit ball is a
perfect dichotomy: it holds in finite dimension, fails in
infinite dimension, with no middle ground --- geometry alone
reads off the dimension type.
\end{example}

\section{Banach spaces}

\begin{definition}\label{def:b2:nvs:banach}
A \emph{Banach space}\index{Banach space} is a complete normed
space. Examples: every finite-dimensional normed space
(\cref{thm:b2:nvs:finitedim}); $\bigl(C(\intcc{a}{b}),
\norm\cdot_\infty\bigr)$ (\cref{thm:b2:metric:rncomplete});
$\mathcal{L}_c(E, F)$ for $F$ Banach (same proof pattern as for
continuous functions). Non-example: $\bigl(C(\intcc{0}{1}),
\norm\cdot_1\bigr)$ (\cref{exo:b2:nvs:7}).
\end{definition}

\begin{example}[The operator norm of integration]\label{ex:b2:nvs:integrationnorm}
On $\bigl(C(\intcc01), \norm\cdot_\infty\bigr)$, let $T(f)(x) =
\int_0^x f(t)\,\dd t$ (an endomorphism: $T(f)$ is continuous).
Run \cref{met:b2:nvs:opnorm}. Upper bound:
\[
\abs{T(f)(x)} \leq \int_0^x\abs f \leq x\,\norm f_\infty \leq
\norm f_\infty ,
\]
so $\vertiii T \leq 1$. Witness: $f \equiv 1$ gives $T(f)(x) =
x$ and $\norm{T(f)}_\infty = 1 = \norm f_\infty$: attained,
$\vertiii T = 1$. But note $\vertiii{T^2} = \frac12 \neq
\vertiii T^2$: indeed $T^2(f)(x) = \int_0^x(x - t)f(t)\dd t$ has
$\abs{T^2(f)(x)} \leq \frac{x^2}2\norm f_\infty$, attained again
at $f \equiv 1$; and generally $\vertiii{T^n} = \frac1{n!}$ ---
the submultiplicative bound $\vertiii T^n = 1$ is off by a
factorial. This is exactly the phenomenon the iterate trick of
\cref{ch:b2:metric}'s weekend problem converts into global
solvability of linear differential equations.
\end{example}

\begin{theorem}[Absolute convergence in Banach spaces]\label{thm:b2:nvs:absoluteconvergence}
In a Banach space, if $\sum \norm{u_n} < \infty$ then $\sum u_n$
converges, and $\norm{\sum u_n} \leq \sum\norm{u_n}$. (The full
theory of series in normed spaces is \cref{ch:b2:series}.)
\end{theorem}

\begin{proof}
Partial sums $S_N$: for $q > p$, $\norm{S_q - S_p} \leq
\sum_{n=p+1}^{q}\norm{u_n}$, which tends to $0$ (Cauchy criterion
for the real series of norms): $(S_N)$ is Cauchy, hence convergent.
The inequality passes to the limit from the finite triangle
inequality.
\end{proof}

\begin{example}[Matrix exponential, first contact]\label{ex:b2:nvs:matrixexp}
$\mathcal{M}_n(K)$ with any submultiplicative norm ($\vertiii{AB}
\leq \vertiii A \vertiii B$) is Banach (finite dimension). Then, for
every $A$,
\[
\eu^A = \sum_{k=0}^{\infty} \frac{A^k}{k!}
\]
converges absolutely ($\vertiii{A^k/k!} \leq \vertiii A^k /k!$,
summable): well defined. \cref{ch:b2:diffeq} exploits it
systematically.
\end{example}

\begin{example}[A Neumann series that terminates]\label{ex:b2:nvs:neumannfinite}
For $A = \begin{psmallmatrix}0 & \frac12\\ 0 &
0\end{psmallmatrix}$: $\vertiii A < 1$ in any operator norm
built on \cref{ex:b2:nvs:examples}'s norms, and $A^2 = 0$, so
the geometric series collapses:
\[
(I - A)^{-1} = \sum_{k \geq 0} A^k = I + A =
\begin{pmatrix}1 & \tfrac12\\ 0 & 1\end{pmatrix},
\]
verified by $(I - A)(I + A) = I - A^2 = I$. Nilpotency
truncates the series exactly as it truncated the exponential in
\cref{ch:b2:reduction}; and the example calibrates
expectations: the Neumann inverse is an infinite series in
general, a polynomial precisely when the perturbation is
nilpotent, and the error after $N$ terms is always bounded by
the geometric tail $\vertiii A^{N+1}/(1 - \vertiii A)$.
\end{example}

\begin{example}[The exponential of a rotation generator]\label{ex:b2:nvs:rotationexp}
Take $A = \begin{psmallmatrix}0 & -\theta\\ \theta &
0\end{psmallmatrix}$. Then $A^2 = -\theta^2 I$, so the powers
cycle with period four, and the series splits into even and odd
parts:
\[
\eu^{A} = \sum_{k}\frac{A^k}{k!}
= \Bigl(\sum_{j}\frac{(-1)^j\theta^{2j}}{(2j)!}\Bigr) I
+ \Bigl(\sum_{j}\frac{(-1)^j\theta^{2j+1}}{(2j+1)!}\Bigr)
\frac{A}{\theta}
= \begin{pmatrix}
\cos\theta & -\sin\theta\\
\sin\theta & \cos\theta
\end{pmatrix},
\]
all rearrangements being licensed by absolute convergence. The
exponential of an antisymmetric generator is a rotation ---
computed here purely from the series, three chapters before the
differential equation $x' = Ax$ (\cref{ch:b2:diffeq}) explains
\emph{why}: $\eu^{tA}$ is uniform circular motion. The closing
insight: identities among matrix series are proved exactly like
scalar ones, once a submultiplicative norm certifies absolute
convergence.
\end{example}

\begin{example}[Sup norm means uniform: the dictionary]\label{ex:b2:nvs:uniformdictionary}
The statement $\norm{f_n - f}_\infty \to 0$ \emph{is} uniform
convergence: one number, $\sup_x\abs{f_n(x) - f(x)}$, bounds
the error at every point simultaneously. The dictionary in
action on $f_n(x) = x^n$ over $\intcc{0}{1}$: pointwise, $f_n
\to 0$ on $\intco{0}{1}$ and $f_n(1) = 1$; in norm,
$\norm{f_n - 0}_\infty = 1 \not\to 0$, and indeed the pointwise
limit is discontinuous, hence out of reach of a
$\norm\cdot_\infty$-limit in $C(\intcc01)$ (which is closed
under uniform limits, \cref{thm:b2:metric:rncomplete}). On
$\intcc{0}{a}$, $a < 1$: $\norm{f_n}_\infty = a^n \to 0$ ---
uniform convergence restored by shrinking the domain. Every
convergence statement of \cref{ch:b2:funcseq} is a statement
about this one norm; keeping the dictionary in mind halves
that chapter.
\end{example}

\begin{remark}[Common pitfalls]\label{rem:b2:nvs:pitfalls}
(i) An operator norm depends on \emph{both} chosen norms: the
same matrix has $\vertiii\cdot_\infty$ given by row sums
(\cref{exo:b2:nvs:4}) and a different $\vertiii\cdot_1$ (column
sums); quoting ``the'' norm of a matrix without naming the
underlying norms is meaningless. (ii) $\vertiii{AB} \leq
\vertiii A\,\vertiii B$ is an inequality, usually strict ---
powers can shrink much faster than the bound $\vertiii A^k$
suggests, which is the whole point of adapted norms (this
chapter's weekend problem, question 22). (iii) ``Linear implies
continuous'' is a finite-dimensional privilege: differentiation
on polynomials is linear and unbounded
(\cref{ex:b2:nvs:operatornorms}). (iv) Absolute convergence of
$\sum u_n$ only helps when the space is complete
(\cref{exo:b2:series:9} builds the counterexample). (v) In
infinite dimension a supremum over the unit ball is a genuine
supremum: do not assume it is attained
(\cref{exo:b2:nvs:8}).
\end{remark}

\begin{remark}[Perspectives within this volume]\label{rem:b2:nvs:perspectives}
Three appointments are now fixed. With \cref{ch:b2:series}: in
a Banach space, absolutely convergent series converge, so the
geometric and exponential series of operators become everyday
tools --- inverting $I - A$, defining $\eu^{A}$
(\cref{ex:b2:series:neumann}). With
\cref{ch:b2:funcseq} and \cref{ch:b2:powerseries}: convergence
of function sequences and power series is convergence in
$\bigl(C, \norm\cdot_\infty\bigr)$
(\cref{ex:b2:nvs:uniformdictionary}), and the radius of
convergence is a statement about which geometric series
dominate. With \cref{ch:b2:fourier}: the norms
$\norm\cdot_2$ and $\norm\cdot_\infty$ genuinely disagree on
$C(\intcc{0}{1})$ (\cref{ex:b2:nvs:sqrtnxn}), which is exactly
why mean-square convergence of Fourier series and uniform
convergence are two different theorems with two different
prices.
\end{remark}

\begin{remark}[Where this chapter is used]
Operator norms and the geometric series drive the perturbation
arguments of \cref{ch:b2:diffcalc} (inverse function theorem) and
the matrix exponential of \cref{ch:b2:diffeq}; equivalence of
norms silently authorizes every ``choose your favourite norm''
argument in \cref{ch:b2:funcseq} and beyond; and the
finite/infinite divide of Riesz's theorem --- made quantitative in
this chapter's weekend problem --- is the reason the Year 3
volume needs new tools (weak convergence, Arzel\`a--Ascoli,
Hilbert space projections) where this volume could still extract
convergent subsequences.
\end{remark}

\section{Exercises}

\begin{exercise}[$\star$]\label{exo:b2:nvs:1}
On $\R^2$, draw the unit balls of $\norm\cdot_1$, $\norm\cdot_2$,
$\norm\cdot_\infty$, and prove the inequalities $\norm x_\infty
\leq \norm x_2 \leq \norm x_1 \leq 2\norm x_\infty$ with the best
constants in dimension $2$.
\end{exercise}

\begin{exercise}[$\star$]\label{exo:b2:nvs:2}
Is $N(f) = \abs{f(0)} + \norm{f'}_\infty$ a norm on
$C^1(\intcc{0}{1})$? Compare it with $\norm{f}_\infty$: one
inequality holds, the other fails (exhibit).
\end{exercise}

\begin{exercise}[$\star$]\label{exo:b2:nvs:3}
Compute the operator norm of $u(f) = \int_0^1 f(t)\,\eu^t\,\dd t$
on $\bigl(C(\intcc{0}{1}), \norm\cdot_\infty\bigr) \to \R$, and of
the shift $S(x_1, x_2, \dots, x_n) = (x_2, \dots, x_n, 0)$ on
$(K^n, \norm\cdot_\infty)$.
\end{exercise}

\begin{exercise}[$\star\star$]\label{exo:b2:nvs:4}
On $(\R^n, \norm\cdot_\infty)$, prove that the operator norm of a
matrix $A$ is $\vertiii A_\infty = \max_i \sum_j \abs{a_{ij}}$ (the
largest absolute row sum). Compute it for $\begin{pmatrix} 1 & -2\\
3 & 1\end{pmatrix}$.
\end{exercise}

\begin{exercise}[$\star\star$]\label{exo:b2:nvs:5}
Prove that $GL_n(K)$ is open in $\mathcal{M}_n(K)$ and that $A
\mapsto A^{-1}$ is continuous on it. \emph{Hint: for openness, if
$\vertiii H < \frac{1}{\vertiii{A^{-1}}}$ then $A + H = A(I +
A^{-1}H)$ with $\vertiii{A^{-1}H} < 1$, and $I + B$ is invertible
for $\vertiii B < 1$ by the geometric series
(\cref{thm:b2:nvs:absoluteconvergence}); for continuity, bound
$(A+H)^{-1} - A^{-1}$ using the same series.}
\end{exercise}

\begin{exercise}[$\star\star$]\label{exo:b2:nvs:6}
Let $\varphi$ be a linear form on a normed space $E$. Prove that
$\varphi$ is continuous if and only if $\ker\varphi$ is closed.
\emph{(If $\ker\varphi$ is closed and $\varphi \neq 0$, pick $a$
with $\varphi(a) = 1$ and $r > 0$ with $B(a, r) \cap \ker\varphi =
\emptyset$; deduce $\abs{\varphi(h)} \leq \frac{1}{r}\norm h$ by a
scaling argument on $a - \frac{h}{\varphi(h)}$.)}
\end{exercise}

\begin{exercise}[$\star\star$]\label{exo:b2:nvs:7}
Prove that $\bigl(C(\intcc{0}{1}), \norm\cdot_1\bigr)$ is not
complete: show that the functions $f_n$, affine ramps from $0$ to
$1$ over $\bigl[\frac12 - \frac1n, \frac12\bigr]$ (value $0$
before, $1$ after), form a Cauchy sequence with no continuous
$\norm\cdot_1$-limit.
\end{exercise}

\begin{exercise}[$\star\star\star$]\label{exo:b2:nvs:8}
On $E = C(\intcc{0}{1})$ with $\norm\cdot_\infty$, consider
\[
\varphi(f) = \sum_{n \geq 1} (-1)^n\, 2^{-n} f\bigl(\tfrac1n\bigr).
\]
Prove that $\varphi$ is a well-defined continuous linear form with
$\vertiii\varphi = 1$, but that the supremum defining
$\vertiii\varphi$ is \emph{not attained} on the closed unit ball.
\emph{(Upper bound: triangle inequality. Norm $= 1$: build
continuous $f_K$ with $\norm{f_K}_\infty \leq 1$ and $f_K(\frac1n)
= (-1)^n$ for $n \leq K$ --- the points $\frac1n$ are isolated from
each other. Non-attainment: equality would force $f(\frac1n) =
(-1)^n$ for every $n$, incompatible with continuity of $f$ at $0$
since $\frac1n \to 0$.)}
\end{exercise}

\begin{exercise}[$\star\star\star$]\label{exo:b2:nvs:9}
Let $E$ be a normed space in which the closed unit ball is compact.
Re-derive, without quoting \cref{thm:b2:nvs:riesz}, that every
bounded sequence has a convergent subsequence, and prove that every
linear form on $E$ is continuous if and only if $\dim E < \infty$.
\emph{(For infinite dimension, build a discontinuous form by
defining it freely on a linearly independent normalized sequence
and extending --- admitting the existence of an algebraic
complement.)}
\end{exercise}

\begin{exercise}[$\star\star$]\label{exo:b2:nvs:10}
On $C(\intcc{0}{1})$, prove $\norm f_1 \leq \norm f_2 \leq \norm
f_\infty$ \emph{(Cauchy--Schwarz for the first)}, and show with
the family $f_n(x) = x^n$ that neither inequality can be reversed
up to a constant: the three norms are pairwise non-equivalent.
\end{exercise}

\begin{exercise}[$\star\star$]\label{exo:b2:nvs:11}
(Distance to a hyperplane) Let $\varphi$ be a nonzero continuous
linear form on a normed space $E$. Prove that
\[
d\bigl(x, \ker\varphi\bigr) =
\frac{\abs{\varphi(x)}}{\vertiii\varphi}
\qquad (x \in E),
\]
and check on \cref{exo:b2:nvs:8} that the infimum need not be
attained by any point of the hyperplane.
\end{exercise}

\begin{exercise}[$\star\star\star$]\label{exo:b2:nvs:12}
On $E = \R[X]$ (all polynomials), let $N_1(P) =
\sup_{\intcc{0}{1}}\abs P$ and $N_2(P) =
\sup_{\intcc{0}{2}}\abs P$. Show that $N_1 \leq N_2$ but that
$N_1$ and $N_2$ are \emph{not} equivalent; deduce that the
identity $(E, N_2) \to (E, N_1)$ is a continuous linear bijection
whose inverse is discontinuous. Show finally that $(E, N_1)$ is
not complete \emph{(Taylor partial sums of $\eu^x$)}. All three
phenomena are impossible in finite dimension --- say why.
\end{exercise}

\section{Problem: Best Approximation and Chebyshev's Theorem}

How well can a function be approximated by polynomials of a given
degree, and which polynomial does it best? On the existence side,
the answer belongs to this chapter: compactness in finite
dimension makes best approximations exist. On the explicit side,
one nontrivial case can be solved completely with bare hands ---
among all \emph{monic} polynomials of degree $n$, the one of
smallest sup norm on $\intcc{-1}{1}$ is the (normalized) Chebyshev
polynomial, of norm $2^{1-n}$: \emph{Chebyshev's extremal
theorem}. The problem proves both sides, then measures how badly
compactness fails in infinite dimension: the unit ball of
$C(\intcc{0}{1})$ contains infinite constellations of points at
mutual distance $1$.

\begin{problem}[{Weekend problem --- Chebyshev's extremal theorem
and the geometry of the unit ball}]
\label{pb:b2:nvs:1}
Norms without subscript are sup norms on the indicated segment.

\medskip
\noindent\textbf{Part I --- Best approximation in normed
spaces.}
\begin{enumerate}
  \item Let $F$ be a finite-dimensional subspace of a normed
        space $E$ and $x \in E$. Prove that the distance $d(x,
        F) = \inf_{f \in F}\norm{x - f}$ is \emph{attained}
        \emph{(reduce to a closed bounded subset of $F$ and use
        \cref{thm:b2:nvs:finitedim})}.
  \item A norm is \emph{strictly convex} when $\norm u = \norm v
        = 1$ and $u \neq v$ imply $\bigl\Vert\frac{u +
        v}2\bigr\Vert < 1$. Show that $\norm\cdot_2$ on $\R^n$
        is strictly convex \emph{(parallelogram identity)}, and
        that $\norm\cdot_1$ and $\norm\cdot_\infty$ are not for
        $n \geq 2$.
  \item Prove that for a strictly convex norm, the best
        approximation of question 1 is \emph{unique}.
  \item In $(\R^2, \norm\cdot_\infty)$, compute all best
        approximations of $x = (0, 1)$ by the line $F =
        \operatorname{Vect}\bigl((1,0)\bigr)$: an interval of
        minimizers.
  \item In $\bigl(C(\intcc{a}{b}), \norm\cdot_\infty\bigr)$,
        show that the best approximation of $f$ by
        \emph{constants} is unique, equal to $c^* = \frac{\max f
        + \min f}{2}$, with distance $\frac{\max f - \min
        f}{2}$; compute both for $f(x) = x^2$ on
        $\intcc{0}{1}$.
\end{enumerate}

\medskip
\noindent\textbf{Part II --- Chebyshev polynomials.}
\begin{enumerate}[resume]
  \item Show there is exactly one polynomial $T_n$ with
        $T_n(\cos\theta) = \cos n\theta$ for all $\theta$
        \emph{(recurrence $T_{n+1} = 2XT_n - T_{n-1}$ from the
        cosine addition formula)}, that $\deg T_n = n$, and that
        its leading coefficient is $2^{n-1}$ for $n \geq 1$.
  \item Show $\abs{T_n} \leq 1$ on $\intcc{-1}{1}$, with
        $T_n(\eta_k) = (-1)^k$ at the $n + 1$ points $\eta_k =
        \cos\frac{k\pi}{n}$ ($k = 0, \dots, n$), and that the
        roots of $T_n$ are the $n$ points
        $\cos\frac{(2k-1)\pi}{2n}$, interlacing the $\eta_k$.
  \item Compute $T_2, T_3, T_4$, and verify the alternation of
        $T_3$ at $\eta_0, \dots, \eta_3 = 1, \frac12, -\frac12,
        -1$ by direct evaluation.
  \item For $\abs x \geq 1$, prove
        \[
        T_n(x) = \frac{\bigl(x + \sqrt{x^2 - 1}\bigr)^n +
        \bigl(x - \sqrt{x^2 - 1}\bigr)^n}{2},
        \]
        and deduce $T_n(x) \sim \frac12\bigl(x + \sqrt{x^2 -
        1}\bigr)^n \to \infty$ geometrically for fixed $x > 1$.
  \item Prove the composition law $T_m \circ T_n = T_{mn}$
        \emph{(check on $\intcc{-1}{1}$ and invoke the rigidity
        of polynomials)}.
\end{enumerate}

\medskip
\noindent\textbf{Part III --- Chebyshev's extremal theorem.}
Write $Q_n = 2^{1-n}T_n$ (monic, by question 6).
\begin{enumerate}[resume]
  \item Let $P$ be monic of degree $n \geq 1$ with
        $\sup_{\intcc{-1}{1}}\abs P < 2^{1-n}$. By evaluating $D
        = Q_n - P$ at the points $\eta_k$ and counting sign
        changes, derive a contradiction. Conclude:
        \[
        \sup_{\intcc{-1}{1}}\abs P \;\geq\; 2^{1-n}
        \qquad\text{for every monic } P \text{ of degree } n.
        \]
  \item (Equality case) Suppose $\sup_{\intcc{-1}{1}}\abs P =
        2^{1-n}$ with $P$ monic of degree $n$, and let $D = Q_n
        - P \neq 0$. Show $(-1)^kD(\eta_k) \geq 0$ for all $k$;
        show that each of the $n$ intervals
        $\intcc{\eta_{k}}{\eta_{k-1}}$ contains a zero of $D$,
        and that a zero shared by two consecutive intervals is
        an \emph{interior} point $\eta_k$ where $D' = 0$ as
        well. Conclude that $D$ has $n$ zeros counted with
        multiplicity, hence $D = 0$: the minimizer is exactly
        $Q_n$ --- \emph{Chebyshev's extremal theorem}.
  \item Restate the theorem as a distance: on $\intcc{-1}{1}$,
        \[
        d_\infty\bigl(X^n,\ \R_{n-1}[X]\bigr) = 2^{1-n},
        \]
        with unique best approximation $X^n - Q_n$; and show by
        the affine substitution $x = \frac{1+t}2$ that on
        $\intcc{0}{1}$ the distance becomes $2^{1-2n}$.
  \item (Optimal interpolation nodes) For $n$ nodes $x_1, \dots,
        x_n \in \intcc{-1}{1}$, the node polynomial
        $\omega(x) = \prod_i(x - x_i)$ is monic of degree $n$.
        Deduce from question 12 which choice of nodes minimizes
        $\sup_{\intcc{-1}{1}}\abs\omega$, the node-dependent
        factor of the classical interpolation error bound, and
        give the minimal value.
  \item Verify the case $n = 2$ of the theorem by hand (find
        $\inf_c \sup_{\intcc{-1}{1}}\abs{x^2 - c}$ directly),
        and compute numerically the distance of question 13 on
        $\intcc{0}{1}$ for $n = 10$. What does its size say
        about the graph of $x^{10}$?
\end{enumerate}

\medskip
\noindent\textbf{Part IV --- The unit ball of
$C(\intcc{0}{1})$.}
\begin{enumerate}[resume]
  \item Let $g_k(x) = x^{2^k}$. Show $\norm{g_k}_\infty = 1$ and
        $\norm{g_k - g_j}_\infty \geq \frac14$ for $j > k$
        \emph{(evaluate at the point where $x^{2^k} =
        \frac12$)}: an explicit bounded sequence with no
        convergent subsequence --- the closed unit ball is not
        compact, by bare hands.
  \item (Riesz's lemma, sharpened) Let $F$ be a
        \emph{finite-dimensional} proper subspace of a normed
        space $E$. Using question 1, produce a unit vector $x$
        with $d(x, F) = 1$ exactly --- not just $\geq 1 -
        \varepsilon$ as in \cref{thm:b2:nvs:riesz}'s lemma.
  \item Deduce: in every infinite-dimensional normed space there
        is a sequence of unit vectors with pairwise distances
        $\geq 1$, and re-derive Riesz's theorem from it.
  \item In $C(\intcc{0}{1})$, exhibit such a constellation
        explicitly: the tent functions $h_n$ supported on
        $\bigl[\frac1{n+1}, \frac1n\bigr]$ with peak value $1$.
        Verify $\norm{h_n} = 1$, $\norm{h_n - h_m} = 1$ for $n
        \neq m$, and note that $h_n \to 0$ pointwise but not
        uniformly.
  \item (Total boundedness fails) Show that the closed unit ball
        of $C(\intcc{0}{1})$ cannot be covered by finitely many
        balls of radius $\frac13$ \emph{(each such ball contains
        at most one $h_n$)} --- contrast with the
        total-boundedness step in the proof of
        \cref{thm:b2:metric:borellebesgue}.
\end{enumerate}

\medskip
\noindent\textbf{Part V --- Norms at work on matrices, and
synthesis.}
\begin{enumerate}[resume]
  \item Prove that every eigenvalue $\lambda$ of $A \in
        \mathcal{M}_n(\C)$ satisfies $\abs\lambda \leq
        \vertiii A$ for every operator norm; apply
        \cref{exo:b2:nvs:4} to bound the eigenvalues of
        $\begin{psmallmatrix}1 & -2\\ 3 & 1\end{psmallmatrix}$
        and compare with their true modulus.
  \item (Adapted norms) Let $A$ be diagonalizable, $A =
        P\,\mathrm{diag}(\lambda_1, \dots, \lambda_n)\,P^{-1}$.
        Show $N_P(x) = \norm{P^{-1}x}_\infty$ is a norm whose
        operator norm satisfies $\vertiii A_{N_P} =
        \max_i\abs{\lambda_i}$.
  \item Deduce: for diagonalizable $A$, $A^k \to 0$ if and only
        if all eigenvalues satisfy $\abs{\lambda_i} < 1$ ---
        equivalence of norms makes the conclusion
        norm-independent. Check on $A =
        \frac14\begin{psmallmatrix}1 & 2\\ 2 &
        1\end{psmallmatrix}$.
  \item (Equivalence constants blow up) On $\R_n[X]$, compare
        $N_c(P) = \max_k \abs{a_k}$ (coefficients) and
        $\norm{P}_{\intcc{0}{1}}$: both are norms, so they are
        equivalent for each fixed $n$; but show, using the monic
        minimizer of question 13 on $\intcc{0}{1}$, that the
        best constant $C_n$ in $N_c \leq
        C_n\norm\cdot_{\intcc{0}{1}}$ satisfies $C_n \geq
        2^{2n-1}$. Conclude in one sentence why ``all norms are
        equivalent'' dies in infinite dimension.
  \item (Synthesis) One sentence each: where compactness of
        finite-dimensional balls worked (questions 1, 12); what
        strict convexity governs; what the constellation of
        questions 18--19 destroys; and how question 24
        quantifies the failure. Name the summit (Chebyshev's
        extremal theorem) and state where best approximation
        finds its modern home (the projection theorem on Hilbert
        spaces, Year 3 volume, where completeness replaces
        compactness).
\end{enumerate}
\end{problem}
